\scanpage{112}
\uline{Topos modulaire pour des structures algébriques
\deleted{\unsure{définies}} sur objets \inserted{fixés} d'un topos fixé \deleted{\unsure{E}} $X$}

\deleted{\unsure{On}} On peut se donner des objets $E_i\in X$, $i\in I$
et \deleted{de s'inté} une espèce de structure $T$ avec
topos modulaire $B$, \deleted{sur objets} avec donnée
``objets de base'' correspondants : des $b_i\in B$, $i\in I$

On s'intéresse donc, pour tt topos $X'$ sur $X$,
\inserted{(ayant comme images inverses des $E_i$ les $E'_i$)}
aux structures d'espèce $T$ \deleted{dans} $X'$ \deleted{qui}
(i.e. morphismes $f:X'\longrightarrow B$) \deleted{tels que} avec
isom donnés $f^*(b_i)\simeq E'_i$. La catégorie
\diagram{\begin{tikzcd}[ampersand replacement=\&,row sep=small]
 B\arrow[d,"b"]\&X\arrow[dl,"E"]\\
 \widehat{(\mathrm{Ens\,f})^o}{}^{I}=B_I\&{}
\end{tikzcd}}
de ces objets sur un topos $X'$
fixé, pour $X'\longrightarrow X$ \uline{et} $X'\longrightarrow B$
variable, est le produit 2-fibré de $\underline{\mathrm{Hom}}_{\mathrm{top}}(X',X)$ et
$\underline{\mathrm{Hom}\,\mathrm{top}}(X',B)$ sur $\underline{\mathrm{Hom}\,\mathrm{top}}(X',B_I)$
Donc on cherche un produit
2-fibré $B\times_{B_I}X$, dans la 2-catégorie (Top)
(On n'a pas encore prouvé que cela marche)
\deleted{Cela doit donc être une sous-\unsure{espèce}}

Supposons qu'il s'agisse d'une $T$ d'une
structure définie par $\varprojlim$ finies, avec cat. modulaire
$R$ (donc $B=\widehat R$ \unsure{$\}$}, $I\lhook\joinrel\longrightarrow\mathrm{Ob}\,R$) Donc on
s'intéresse aux foncteurs exact à g.
$R\longrightarrow X'$ transformant les $r_i$ en les $E'_i$.
Soit $\unsure{\mathscr S|S|\mathcal S}_{\deleted{\unsure{X}}}$ \deleted{la catégorie}\inserted{le champ sur $X$} formée des structures
de ce type dans $(X,(E_i))$, \inserted{et les fonctrs induits} i.e. des functs ex
à g. $f:R\longrightarrow X$ avec $f(r_i)\simeq E_i$ donnés, et
% slanted
\margin{champ en \unsure{cat.} à choisir (associé à un faisc. en cat. donné ?)}
\deleted{et soit Considérons la}
soit. $\unsure{\Phi|I|\underline I}_X$ le champ \inserted{en topos} des flèches de $\deleted{\unsure{B}}X$.
% archive sheet number: 111
